\scanpage{61}

qui généralise (126).

\deleted{Quand on part d'une situation \unsure{isomorphe} à un}

Pour appliquer ces réflexions au cas d'un $M$ \deleted{\unsure{libre sur|libre}} de rang 2 sur un anneau $\Lambda_n=\mathbb Z/n\mathbb Z$, il faudrait vérifier que pour

\[
 i=\{e,-e\}\in(\text{dét}\,M)^{*}/\{\pm1_\Lambda\}
\]

d'où un $A_{M,i}$, \deleted{\unsure{les élts de son}} \inserted{$\simeq\deleted{\unsure{R}}\operatorname{Rep}(M)/\pm1\subset\unsure{\mathscr P|\mathcal P|\mathfrak P}_2(S_{\unsure{M}})$} où $S=S_M=M^{*}/\{\pm\mathrm{id}_M\}$, on a \deleted{déjà}

\[
 \underset{\substack{\cap\\ \unsure{\mathfrak S|\mathscr S}_S}}{Z_0}
 =\deleted{\unsure{Z'_n\ \text{opérant par}}}
 \underbrace{\{T_\lambda\mid\lambda\in{Z'_n}=\Lambda^*/\pm1_\Lambda\}}_{\text{noté }(Z'_n)_S},
 \tag*{(154)}
\]

et on suspecte même que l'on a déjà

\[
 \operatorname{Comm}_{\mathfrak S_S}
 \bigl(\underset{\substack{\Vert\\{\mathrm{Gl}}^{\unsure{1|\prime}}_\Lambda(M)^\pm}}{G_0}\bigr)
 \overset{?}{=}(Z'_n)_S\ \unsure{;|.}
 \tag*{(155)}
\]

Mais sans s'encombrer de l'hypothèse $\Lambda\simeq\mathbb Z/n\mathbb Z$, on peut regarder sur $S_M=\underbrace{M^{**}}_{\substack{\text{ens des}\\\text{él. de }M\\\text{qui font partie d'une}\\\text{base}}}/\{\pm\mathrm{id}_M\}$ la structure de graphe définie par $A_{M,i}$ $(i=\{e,-e\}\in(\text{det}\,M)^*/\pm1_\Lambda))$ et les groupes $G_0,Z_0$ associés.

% archive sheet number: 60
